% !TEX root = build_experiments.tex
% Section 5 (Conclusion with limitations and future work), v2 2026-11-08, written to the user's
% positioning: a baseline toward collision-free execution of frozen VLAs without an external segmenter;
% future work on two complementary axes (perception: synchronizing the policy's internal representation
% with scene geometry; action: constraint-aware chunk generation, ongoing).
% Red = depends on projected results or is a proposed statement; black = method facts and cited work.
\section{Conclusion}
\label{sec:conclusion}

We presented \sys, which builds the collision constraints a frozen VLA needs from the policy's own representations and the observed scene geometry, without an external segmenter and without changing the policy. Attention names the target, depth continuity bounds its observed surface, a spatiotemporal memory keeps its identity, and a stage- and link-conditioned contact policy decides which link may enter it and when; a fixed trajectory optimizer enforces the result. \proj{On SafeLIBERO and LIBERO-Safety with the same frozen $\pi_{0.5}$, this roughly halves the collisions of the strongest safety layer and gives the highest safe-success rate, with task success on par with the best layer and 1.4 points below the policy alone on obstacle-free LIBERO, and the ablations show that the gain requires both the depth-bounded surface and the task-conditioned policy.} We see \sys\ as a baseline for collision-aware execution of frozen VLAs: a segmenter-free, policy-grounded collision representation whose remaining failures are attributable.

\textbf{Limitations and future work.} \proj{Two limitations follow from the failure analysis of Sec.~\ref{sec:exp:analysis}, and they are complementary. On the perception side, \sys\ reads geometry out of a policy that was never trained to represent it: the target surface is only as complete as the observation, and the attention and the identity features can drift from the scene under occlusion or clutter. Frozen VLAs already encode more about the task than their actions reveal~\cite{zhang2026frozenvlas,bhardwaj2026progress}; the next step is to bring this internal representation and the observed geometry closer together, through explicit 3D representations in the backbone~\cite{qu2025spatialvla,li2025pointvla,liu2026activevla} or training signals that make the transformer's features geometry-aware, so that the same read-out becomes more reliable. On the action side, the trajectory optimizer refines a chunk after it is generated: it can deflect a pose but cannot produce a new one, so a corridor narrower than the margin ends in a hold, and a deflected chunk can disturb the grasp or placement pose. The same separate refinement pass is also where the cost lies: the surface, the fields, and the optimization run once per chunk and leave \sys\ one to two orders of magnitude slower per request than an ellipsoid filter, which we report as a limitation rather than claim real-time operation. Resolving both requires the constraints to shape the generation of the action chunk itself, including the 6D grasp and placement poses, in the spirit of constrained or guided generation~\cite{english2026constrainedfm,song2026omniguide,jia2026guidedvla}; we are pursuing this in ongoing work. Both directions build on the representation introduced here.}
